\chapter{Introduction}
\label{ch:intro}

\section{Context}
\label{sec:intro_contexte}

The \esibot{} project is part of an IoT / Mobile Robotics integrative module at the
École Nationale Supérieure d'Informatique (ESI Algiers). The objective of this project
is to design, build and program an autonomous mobile robot capable of mapping an unknown
environment and navigating within it to defined goal points. This project is inspired by
the architecture of the \textbf{TurtleBot3 Burger}, replacing its proprietary hardware
components with accessible, low-cost solutions: an ESP32-CAM \gls{mcu} microcontroller, a
Raspberry Pi 4, a Kinect depth sensor, and 3D printing for the chassis.

\section{Report Objectives}
\label{sec:intro_objectifs}

This document details the hardware, software and mechanical architecture of the system.
It highlights the final design choices that led to the robot's implementation, in
particular the integration of motor control on the Raspberry\,Pi\,4 and the use of the
Kinect sensor for mapping.

The report covers:
\begin{itemize}
	\item The complete architecture (mechanical, electronic, and software components);
	\item Sensor interfacing and communication flows under \ros{};
	\item Configuration parameters and sensor fusion;
	\item Performance evaluation and the system's limitations.
\end{itemize}

\section{Report Organization}
\label{sec:intro_organisation}

This report is structured as follows:
Chapter~\ref{ch:archi_globale} details the overall architecture. Chapter~\ref{ch:mecanique}
covers the mechanical and electronic design. Chapter~\ref{ch:commande_bas_niveau} documents
low-level control and data fusion. The mapping, navigation and exploration approach is
covered in Chapter~\ref{ch:slam}. Chapter~\ref{ch:ihm} presents the user interface and
deployment procedures. Finally, the integration results, limitations analysis, and project
outlook are addressed in the last three chapters
(\ref{ch:resultats}, \ref{ch:limitations}, \ref{ch:bilan}).
